\documentclass[11pt,a4paper]{article}

\usepackage[margin=1in]{geometry}
\usepackage{amsmath,amssymb,amsthm}
\usepackage{mathtools}
\usepackage{booktabs}
\usepackage{enumitem}
\usepackage{fancyhdr}
\usepackage{titlesec}
\usepackage{hyperref}
\usepackage{xcolor}

\hypersetup{
    colorlinks=true,
    linkcolor=blue!70!black,
    citecolor=green!50!black,
    urlcolor=blue!60!black,
}

\newtheorem{theorem}{Theorem}[section]
\newtheorem{lemma}[theorem]{Lemma}
\newtheorem{proposition}[theorem]{Proposition}
\newtheorem{corollary}[theorem]{Corollary}
\theoremstyle{definition}
\newtheorem{definition}[theorem]{Definition}
\theoremstyle{remark}
\newtheorem{remark}[theorem]{Remark}

\DeclareMathOperator*{\argmin}{arg\,min}
\DeclareMathOperator*{\argmax}{arg\,max}
\DeclareMathOperator{\sign}{sign}
\DeclareMathOperator{\Var}{Var}
\DeclareMathOperator{\diag}{diag}
\DeclareMathOperator{\softmax}{softmax}
\DeclareMathOperator{\tr}{tr}

\pagestyle{fancy}
\fancyhf{}
\fancyhead[L]{\textit{scry-learn: Mathematical Foundations}}
\fancyhead[R]{\thepage}
\renewcommand{\headrulewidth}{0.4pt}

\titleformat{\section}{\Large\bfseries}{}{0em}{}[\titlerule]
\titleformat{\subsection}{\large\bfseries}{}{0em}{}

\title{\vspace{-1cm}\textbf{scry-learn: Mathematical Foundations}\\[0.3em]
\large Complete Derivations and Proofs for All ML Algorithms}
\author{scry-learn Technical Reference}
\date{}

\begin{document}
\maketitle
\thispagestyle{fancy}

\noindent\textbf{Notation.} Vectors are bold lowercase ($\mathbf{x}$), matrices bold uppercase ($\mathbf{X}$), scalars italic ($\alpha$). Data is column-major: \texttt{features[feature\_idx][sample\_idx]}. We use $n$ for sample count, $m$ for feature count, $K$ for class count. $\mathbf{1}$ denotes the all-ones vector.

\tableofcontents
\newpage

%=============================================================================
\section{Linear Models}
%=============================================================================

\subsection{Ordinary Least Squares (OLS)}
\label{sec:ols}

\noindent\textit{Source:} \texttt{crates/scry-learn/src/linear/regression.rs}

\subsubsection*{Objective}

Minimize the sum of squared residuals:
\begin{equation}
\hat{\boldsymbol{\beta}} = \argmin_{\boldsymbol{\beta}} \frac{1}{2}\|\mathbf{y} - \mathbf{X}\boldsymbol{\beta} - \beta_0\mathbf{1}\|^2
\end{equation}

\subsubsection*{Normal Equations: Derivation}

Setting the gradient to zero:
\begin{align}
\nabla_{\boldsymbol{\beta}} \mathcal{L} &= -\mathbf{X}^\top(\mathbf{y} - \mathbf{X}\boldsymbol{\beta}) = \mathbf{0} \\
\Rightarrow \quad \mathbf{X}^\top\mathbf{X}\boldsymbol{\beta} &= \mathbf{X}^\top\mathbf{y} \\
\Rightarrow \quad \hat{\boldsymbol{\beta}} &= (\mathbf{X}^\top\mathbf{X})^{-1}\mathbf{X}^\top\mathbf{y}
\end{align}

\begin{theorem}[Optimality of OLS]
The Hessian $\mathbf{H} = \mathbf{X}^\top\mathbf{X}$ is positive semi-definite since $\mathbf{v}^\top\mathbf{X}^\top\mathbf{X}\mathbf{v} = \|\mathbf{X}\mathbf{v}\|^2 \geq 0$ for all $\mathbf{v}$. When $\mathbf{X}$ has full column rank, $\mathbf{H}$ is positive definite and the solution is a unique global minimum.
\end{theorem}

\subsubsection*{Intercept Handling}

The implementation centers the target: $\bar{y} = \frac{1}{n}\sum_i y_i$, then recovers the intercept as $\beta_0 = \bar{y} - \sum_j \beta_j \bar{x}_j$.

\subsubsection*{Solvers}

\textbf{Normal (Gauss-Jordan with partial pivoting).} Solves $(\mathbf{X}^\top\mathbf{X} + \alpha\mathbf{I})\boldsymbol{\beta} = \mathbf{X}^\top\mathbf{y}$ via in-place row reduction. Singularity threshold: $10^{-12}$. Complexity: $O(m^3)$ for the solve step, $O(nm^2)$ to form the Gram matrix.

\textbf{QR Decomposition.} Factor $\mathbf{X} = \mathbf{Q}\mathbf{R}$ via Householder reflections. Then:
\begin{equation}
\hat{\boldsymbol{\beta}} = \mathbf{R}^{-1}\mathbf{Q}^\top\mathbf{y}
\end{equation}
More numerically stable: the condition number is $\kappa(\mathbf{X})$ versus $\kappa(\mathbf{X})^2$ for the normal equations.

\begin{proposition}[QR Stability]
The Householder QR factorization satisfies $\mathbf{Q}^\top\mathbf{Q} = \mathbf{I}$ to machine precision. The triangular solve $\mathbf{R}\boldsymbol{\beta} = \mathbf{Q}^\top\mathbf{y}$ has backward error $O(\kappa(\mathbf{R})\epsilon_{\mathrm{mach}})$, whereas the normal equations have backward error $O(\kappa(\mathbf{X})^2\epsilon_{\mathrm{mach}})$.
\end{proposition}

\textbf{SVD.} Factor $\mathbf{X} = \mathbf{U}\boldsymbol{\Sigma}\mathbf{V}^\top$. The pseudoinverse gives:
\begin{equation}
\hat{\boldsymbol{\beta}} = \mathbf{V}\boldsymbol{\Sigma}^+\mathbf{U}^\top\mathbf{y}, \quad \sigma_i^+ = \begin{cases} 1/\sigma_i & \text{if } \sigma_i > \epsilon \\ 0 & \text{otherwise} \end{cases}
\end{equation}
Handles rank-deficient and wide ($n < m$) matrices. The implementation uses one-sided Jacobi rotation SVD and reports condition number $\kappa = \sigma_{\max}/\sigma_{\min}$ and effective rank.

\textbf{Auto.} Uses Normal equations, falls back to SVD if the matrix is detected as singular.

\subsubsection*{Sparse Support}

When data contains a \texttt{CscMatrix}, the implementation computes $\mathbf{X}^\top\mathbf{X}$ and $\mathbf{X}^\top\mathbf{y}$ directly from compressed sparse column format, iterating only non-zero entries. Complexity: $O(\mathrm{nnz} \cdot m)$ instead of $O(nm^2)$.

%-----------------------------------------------------------------------------
\subsection{Ridge Regression (L2 Regularization)}
\label{sec:ridge}

\noindent\textit{Source:} \texttt{crates/scry-learn/src/linear/regression.rs} (parameter \texttt{alpha > 0})

\subsubsection*{Objective}
\begin{equation}
\hat{\boldsymbol{\beta}} = \argmin_{\boldsymbol{\beta}} \frac{1}{2n}\|\mathbf{y} - \mathbf{X}\boldsymbol{\beta}\|^2 + \frac{\alpha}{2}\|\boldsymbol{\beta}\|_2^2
\end{equation}

\subsubsection*{Closed-Form Solution}
\begin{equation}
\hat{\boldsymbol{\beta}} = (\mathbf{X}^\top\mathbf{X} + n\alpha\mathbf{I})^{-1}\mathbf{X}^\top\mathbf{y}
\end{equation}

\begin{theorem}[Ridge Regularization Guarantees]
For any $\alpha > 0$, the matrix $\mathbf{X}^\top\mathbf{X} + n\alpha\mathbf{I}$ is positive definite (its eigenvalues are $\lambda_i + n\alpha > 0$), guaranteeing a unique solution even when $\mathbf{X}^\top\mathbf{X}$ is singular.
\end{theorem}

\begin{proof}
Let $\lambda_1, \ldots, \lambda_m$ be the eigenvalues of $\mathbf{X}^\top\mathbf{X}$ (all $\geq 0$). The eigenvalues of $\mathbf{X}^\top\mathbf{X} + n\alpha\mathbf{I}$ are $\lambda_i + n\alpha > 0$ for $\alpha > 0$. Therefore the matrix is positive definite and invertible.
\end{proof}

\begin{remark}
Ridge regression trades bias for variance: the solution $\hat{\boldsymbol{\beta}}_\alpha$ satisfies $\|\hat{\boldsymbol{\beta}}_\alpha\|_2 \leq \|\hat{\boldsymbol{\beta}}_{\mathrm{OLS}}\|_2$. This shrinkage reduces variance at the cost of introducing bias, which can improve generalization on high-dimensional or collinear data.
\end{remark}

%-----------------------------------------------------------------------------
\subsection{Lasso Regression (L1 Regularization)}
\label{sec:lasso}

\noindent\textit{Source:} \texttt{crates/scry-learn/src/linear/lasso.rs}

\subsubsection*{Objective}
\begin{equation}
\hat{\boldsymbol{\beta}} = \argmin_{\boldsymbol{\beta}} \frac{1}{2n}\|\mathbf{y} - \mathbf{X}\boldsymbol{\beta} - \beta_0\mathbf{1}\|^2 + \alpha\|\boldsymbol{\beta}\|_1
\end{equation}

\subsubsection*{Coordinate Descent Algorithm}

Since the L1 norm is non-differentiable at zero, we cannot use gradient-based methods directly. Coordinate descent optimizes one coefficient at a time while holding others fixed.

For coordinate $j$, define the partial residual:
\begin{equation}
r_i^{(j)} = y_i - \beta_0 - \sum_{k \neq j} x_{ik}\beta_k
\end{equation}

The coordinate-wise subproblem becomes:
\begin{equation}
\min_{\beta_j} \frac{1}{2n}\sum_{i=1}^n (r_i^{(j)} - x_{ij}\beta_j)^2 + \alpha|\beta_j|
\end{equation}

Computing $\rho_j = \frac{1}{n}\sum_{i=1}^n x_{ij}r_i^{(j)} = \frac{1}{n}\mathbf{x}_j^\top\mathbf{r}^{(j)}$ and defining $z_j = \frac{1}{n}\|\mathbf{x}_j\|^2$, the solution is:

\begin{equation}
\hat{\beta}_j = \frac{S(\rho_j, \alpha)}{z_j}
\end{equation}

where $S$ is the soft-thresholding operator.

\begin{definition}[Soft-Thresholding Operator]
\begin{equation}
S(z, \gamma) = \sign(z)\max(|z| - \gamma, 0) = \begin{cases}
z - \gamma & \text{if } z > \gamma \\
0 & \text{if } |z| \leq \gamma \\
z + \gamma & \text{if } z < -\gamma
\end{cases}
\end{equation}
\end{definition}

\begin{theorem}[Soft-Thresholding is the Proximal Operator of L1]
\label{thm:soft-thresh}
The soft-thresholding operator $S(z, \gamma)$ is the proximal operator of $\gamma|\cdot|$:
\begin{equation}
S(z, \gamma) = \argmin_u \left\{\frac{1}{2}(u - z)^2 + \gamma|u|\right\}
\end{equation}
\end{theorem}

\begin{proof}
Let $f(u) = \frac{1}{2}(u-z)^2 + \gamma|u|$. We consider three cases:

\textbf{Case 1:} $u > 0$. Then $f(u) = \frac{1}{2}(u-z)^2 + \gamma u$. Setting $f'(u) = u - z + \gamma = 0$ gives $u = z - \gamma$. This is valid (i.e., $u > 0$) when $z > \gamma$.

\textbf{Case 2:} $u < 0$. Then $f(u) = \frac{1}{2}(u-z)^2 - \gamma u$. Setting $f'(u) = u - z - \gamma = 0$ gives $u = z + \gamma$. This is valid when $z < -\gamma$.

\textbf{Case 3:} $u = 0$. $f(0) = \frac{1}{2}z^2$. Comparing with the boundary values of Cases 1 and 2 at $u \to 0^+$ and $u \to 0^-$, $u = 0$ is optimal when $|z| \leq \gamma$.

Combining all cases yields the soft-thresholding formula.
\end{proof}

\begin{proposition}[Convergence of Coordinate Descent for Lasso]
Coordinate descent for the Lasso converges to the global optimum. The objective is convex (sum of convex functions) and the L1 penalty is separable across coordinates ($\|\boldsymbol{\beta}\|_1 = \sum_j |\beta_j|$). By Tseng (2001), block coordinate descent converges for convex objectives with separable non-smooth terms.
\end{proposition}

\subsubsection*{Intercept Update}

The intercept is not regularized and is updated each iteration:
\begin{equation}
\beta_0 \leftarrow \frac{1}{n}\sum_{i=1}^n (y_i - \mathbf{x}_i^\top\boldsymbol{\beta})
\end{equation}

%-----------------------------------------------------------------------------
\subsection{Elastic Net (L1 + L2 Regularization)}
\label{sec:elastic-net}

\noindent\textit{Source:} \texttt{crates/scry-learn/src/linear/elastic\_net.rs}

\subsubsection*{Objective}
\begin{equation}
\hat{\boldsymbol{\beta}} = \argmin_{\boldsymbol{\beta}} \frac{1}{2n}\|\mathbf{y} - \mathbf{X}\boldsymbol{\beta}\|^2 + \alpha\cdot\ell_1 \cdot \|\boldsymbol{\beta}\|_1 + \frac{\alpha(1-\ell_1)}{2}\|\boldsymbol{\beta}\|_2^2
\end{equation}
where $\ell_1 \in [0,1]$ is the \texttt{l1\_ratio} parameter interpolating between Ridge ($\ell_1 = 0$) and Lasso ($\ell_1 = 1$).

\subsubsection*{Coordinate Descent Update}

Defining $\rho_j = \frac{1}{n}\mathbf{x}_j^\top\mathbf{r}^{(j)}$, the L1 penalty $\lambda_1 = \alpha \cdot \ell_1$, and the L2 penalty $\lambda_2 = \alpha(1-\ell_1)$:

\begin{equation}
\hat{\beta}_j = \frac{S(\rho_j, \lambda_1)}{\frac{1}{n}\|\mathbf{x}_j\|^2 + \lambda_2}
\end{equation}

\begin{remark}
The L2 term modifies the denominator, providing both shrinkage (via the denominator) and sparsity (via the numerator's soft-thresholding). This combined penalty addresses two limitations: Lasso's instability with correlated features (grouping effect) and Ridge's inability to perform feature selection.
\end{remark}

%-----------------------------------------------------------------------------
\subsection{Logistic Regression}
\label{sec:logreg}

\noindent\textit{Source:} \texttt{crates/scry-learn/src/linear/logistic.rs}

\subsubsection*{Binary Classification}

\textbf{Sigmoid function:}
\begin{equation}
\sigma(z) = \frac{1}{1 + e^{-z}}, \quad \sigma'(z) = \sigma(z)(1 - \sigma(z))
\end{equation}

\textbf{Model:} $P(y=1|\mathbf{x}) = \sigma(\mathbf{w}^\top\mathbf{x} + b)$

\textbf{Negative log-likelihood (binary cross-entropy):}
\begin{equation}
\mathcal{L}(\mathbf{w}) = -\frac{1}{n}\sum_{i=1}^n \left[y_i\log\sigma(z_i) + (1-y_i)\log(1-\sigma(z_i))\right] + \frac{\alpha}{2}\|\mathbf{w}\|_2^2
\end{equation}
where $z_i = \mathbf{w}^\top\mathbf{x}_i + b$.

\textbf{Gradient:}
\begin{equation}
\nabla_\mathbf{w}\mathcal{L} = \frac{1}{n}\mathbf{X}^\top(\boldsymbol{\sigma} - \mathbf{y}) + \alpha\mathbf{w}
\end{equation}
where $\boldsymbol{\sigma} = [\sigma(z_1), \ldots, \sigma(z_n)]^\top$. This uses the identity $\frac{\partial}{\partial z}\log\sigma(z) = 1 - \sigma(z)$ and $\frac{\partial}{\partial z}\log(1-\sigma(z)) = -\sigma(z)$.

\begin{remark}[Binary Fast Path]
For 2-class problems, scry-learn uses a single weight vector with sigmoid instead of 2-class softmax, halving the parameter count from $2m$ to $m$.
\end{remark}

\subsubsection*{Multiclass Classification (Softmax)}

\textbf{Softmax function:}
\begin{equation}
P(y=k|\mathbf{x}) = \frac{\exp(\mathbf{w}_k^\top\mathbf{x} + b_k)}{\sum_{j=1}^K\exp(\mathbf{w}_j^\top\mathbf{x} + b_j)}
\end{equation}

\textbf{Cross-entropy loss:}
\begin{equation}
\mathcal{L} = -\frac{1}{n}\sum_{i=1}^n\sum_{k=1}^K \mathbb{1}[y_i = k]\log P(y_i = k|\mathbf{x}_i) + \frac{\alpha}{2}\sum_{k=1}^K\|\mathbf{w}_k\|_2^2
\end{equation}

\textbf{Gradient for class $k$:}
\begin{equation}
\nabla_{\mathbf{w}_k}\mathcal{L} = \frac{1}{n}\sum_{i=1}^n (P(y=k|\mathbf{x}_i) - \mathbb{1}[y_i = k])\mathbf{x}_i + \alpha\mathbf{w}_k
\end{equation}

\subsubsection*{L1 Proximal Gradient Step}

For L1 or ElasticNet penalties (GD solver only), each gradient step is followed by a proximal step:
\begin{equation}
\mathbf{w} \leftarrow S(\mathbf{w} - \eta\nabla_\mathbf{w}\mathcal{L}_{\text{smooth}}, \eta\lambda_1)
\end{equation}
where $\mathcal{L}_{\text{smooth}}$ is the loss plus L2 term, and $\lambda_1 = \alpha \cdot \ell_1$.

%-----------------------------------------------------------------------------
\subsection{L-BFGS Optimizer}
\label{sec:lbfgs}

\noindent\textit{Source:} \texttt{crates/scry-learn/src/linear/lbfgs.rs}

L-BFGS (Limited-memory Broyden--Fletcher--Goldfarb--Shanno) is the default solver for logistic regression.

\subsubsection*{Two-Loop Recursion}

Given the most recent $p$ correction pairs $\{(\mathbf{s}_i, \mathbf{y}_i)\}_{i=k-p}^{k-1}$ where $\mathbf{s}_i = \mathbf{x}_{i+1} - \mathbf{x}_i$ and $\mathbf{y}_i = \nabla f_{i+1} - \nabla f_i$:

\begin{enumerate}[label=\arabic*., itemsep=0pt]
\item $\mathbf{q} \leftarrow \nabla f_k$
\item For $i = k-1, k-2, \ldots, k-p$: \quad $\alpha_i \leftarrow \frac{\mathbf{s}_i^\top\mathbf{q}}{\mathbf{y}_i^\top\mathbf{s}_i}$; \quad $\mathbf{q} \leftarrow \mathbf{q} - \alpha_i\mathbf{y}_i$
\item $\gamma \leftarrow \frac{\mathbf{s}_{k-1}^\top\mathbf{y}_{k-1}}{\mathbf{y}_{k-1}^\top\mathbf{y}_{k-1}}$; \quad $\mathbf{r} \leftarrow \gamma\mathbf{q}$ \hfill (initial Hessian scaling $H_0 = \gamma I$)
\item For $i = k-p, k-p+1, \ldots, k-1$: \quad $\beta \leftarrow \frac{\mathbf{y}_i^\top\mathbf{r}}{\mathbf{y}_i^\top\mathbf{s}_i}$; \quad $\mathbf{r} \leftarrow \mathbf{r} + (\alpha_i - \beta)\mathbf{s}_i$
\item Return search direction $\mathbf{d} = -\mathbf{r}$
\end{enumerate}

The implementation uses history size $p = 10$.

\subsubsection*{Line Search (Strong Wolfe Conditions)}

Find step size $t > 0$ satisfying:
\begin{align}
f(\mathbf{x}_k + t\mathbf{d}_k) &\leq f(\mathbf{x}_k) + c_1 t \nabla f_k^\top\mathbf{d}_k && \text{(sufficient decrease)} \\
|\nabla f(\mathbf{x}_k + t\mathbf{d}_k)^\top\mathbf{d}_k| &\leq c_2 |\nabla f_k^\top\mathbf{d}_k| && \text{(curvature)}
\end{align}
with $c_1 = 10^{-4}$, $c_2 = 0.9$. Backtracking factor $\rho = 0.5$, maximum 20 backtracks.

\begin{proposition}[L-BFGS Superlinear Convergence]
For strongly convex functions with Lipschitz continuous gradients, L-BFGS converges superlinearly under mild conditions on the correction pairs (Nocedal \& Wright, 2006).
\end{proposition}


%=============================================================================
\section{Decision Trees (CART)}
%=============================================================================

\noindent\textit{Source:} \texttt{crates/scry-learn/src/tree/cart/builder.rs}

\subsection{Split Criteria}

\begin{definition}[Gini Impurity]
For a node with class proportions $p_1, \ldots, p_K$:
\begin{equation}
G = 1 - \sum_{k=1}^K p_k^2
\end{equation}
Measures the probability of misclassifying a randomly chosen sample. $G = 0$ when pure (all same class), $G = 1 - 1/K$ at maximum impurity.
\end{definition}

\begin{definition}[Entropy]
\begin{equation}
H = -\sum_{k=1}^K p_k \log_2 p_k
\end{equation}
From information theory: the expected number of bits to encode the class label. $H = 0$ when pure, $H = \log_2 K$ at maximum entropy.
\end{definition}

\begin{definition}[MSE (Regression)]
\begin{equation}
\text{MSE} = \frac{1}{|S|}\sum_{i \in S}(y_i - \bar{y}_S)^2
\end{equation}
where $\bar{y}_S = \frac{1}{|S|}\sum_{i \in S} y_i$.
\end{definition}

\subsection{Split Selection}

For a candidate split dividing node $S$ into $S_L$ (left) and $S_R$ (right):

\begin{equation}
\Delta I = I(S) - \frac{|S_L|}{|S|}I(S_L) - \frac{|S_R|}{|S|}I(S_R)
\end{equation}

where $I$ is the impurity measure (Gini, Entropy, or MSE). The best split maximizes $\Delta I$ across all features and thresholds.

\subsubsection*{Implementation: Pre-Sorted Indices}

Features are sorted once ($O(nm\log n)$), producing per-feature sorted index arrays. At each node, the split search scans each feature's sorted indices, maintaining running left/right class counts for $O(1)$ impurity updates per candidate threshold. Total split-finding cost per node: $O(nm)$.

\subsection{Cost-Complexity Pruning}

\begin{definition}[Cost-Complexity]
For a subtree $T$, the cost-complexity measure is:
\begin{equation}
R_\alpha(T) = R(T) + \alpha|T|
\end{equation}
where $R(T)$ is the total impurity (weighted by sample count), $|T|$ is the number of leaf nodes, and $\alpha \geq 0$ is the complexity parameter.
\end{definition}

\begin{theorem}[Minimal Cost-Complexity Subtree]
For every $\alpha$, there exists a unique smallest subtree $T(\alpha)$ of the full tree $T_{\max}$ that minimizes $R_\alpha$. The sequence of optimal subtrees forms a nested chain: $T(0) = T_{\max} \supseteq T(\alpha_1) \supseteq T(\alpha_2) \supseteq \cdots \supseteq \{root\}$.
\end{theorem}

The implementation parameter \texttt{ccp\_alpha} controls the pruning strength.

\subsection{FlatTree Optimization}

After building the recursive tree, it is flattened into a contiguous \texttt{Vec<FlatNode>} for cache-optimal prediction. Each \texttt{FlatNode} stores the split feature, threshold, and left/right child indices. Prediction traverses the flat array without pointer chasing.


%=============================================================================
\section{Ensemble Methods}
%=============================================================================

\subsection{Random Forest}
\label{sec:rf}

\noindent\textit{Source:} \texttt{crates/scry-learn/src/tree/random\_forest.rs}

\subsubsection*{Algorithm}

For $b = 1, \ldots, B$ (number of estimators):
\begin{enumerate}[label=\arabic*., itemsep=0pt]
\item Draw a bootstrap sample $\mathcal{D}_b^*$ of size $n$ (with replacement) from $\mathcal{D}$.
\item Grow a CART tree on $\mathcal{D}_b^*$, at each split considering only a random subset of $q$ features.
\end{enumerate}

Prediction: majority vote (classification) or mean (regression) across all $B$ trees.

\subsubsection*{Feature Subsampling}

The number of candidate features per split:
\begin{center}
\begin{tabular}{ll}
\toprule
\texttt{MaxFeatures} & $q$ \\
\midrule
\texttt{Sqrt} (default, classification) & $\lfloor\sqrt{m}\rfloor$ \\
\texttt{Log2} & $\lfloor\log_2 m\rfloor$ \\
\texttt{All} & $m$ \\
\texttt{Fixed}$(k)$ & $k$ \\
\bottomrule
\end{tabular}
\end{center}

\begin{theorem}[Variance Reduction through Bagging]
\label{thm:bag-var}
Let $T_1, \ldots, T_B$ be trees with pairwise correlation $\rho$ and individual variance $\sigma^2$. The variance of the ensemble average is:
\begin{equation}
\Var\left[\frac{1}{B}\sum_{b=1}^B T_b\right] = \rho\sigma^2 + \frac{1-\rho}{B}\sigma^2
\end{equation}
As $B \to \infty$, the variance approaches $\rho\sigma^2$. Feature subsampling reduces $\rho$, further reducing ensemble variance.
\end{theorem}

\subsubsection*{Out-of-Bag (OOB) Error}

Each tree $T_b$ is trained on $\sim 63.2\%$ of the data (the expected fraction of unique samples in a bootstrap of size $n$). The remaining $\sim 36.8\%$ form the OOB set for $T_b$. OOB predictions are aggregated via atomic arrays across parallel tree training (rayon), providing an unbiased estimate of generalization error without a separate validation set.

\begin{proposition}[OOB Approximates Leave-One-Out CV]
The OOB error estimate is approximately equivalent to leave-one-out cross-validation, as each sample is evaluated only by the trees that did not include it in their bootstrap sample.
\end{proposition}

%-----------------------------------------------------------------------------
\subsection{Gradient Boosting}
\label{sec:gbt}

\noindent\textit{Source:} \texttt{crates/scry-learn/src/tree/gradient\_boosting.rs}

\subsubsection*{Functional Gradient Descent Framework}

Initialize $F_0(\mathbf{x}) = \argmin_c \sum_{i=1}^n L(y_i, c)$. For $t = 1, \ldots, T$:

\begin{enumerate}[label=\arabic*., itemsep=0pt]
\item Compute pseudo-residuals: $\tilde{r}_{it} = -\frac{\partial L(y_i, F_{t-1}(\mathbf{x}_i))}{\partial F_{t-1}(\mathbf{x}_i)}$
\item Fit a regression tree $h_t$ to targets $\{\tilde{r}_{it}\}$
\item Update leaf values (terminal region fitting)
\item $F_t(\mathbf{x}) = F_{t-1}(\mathbf{x}) + \eta \cdot h_t(\mathbf{x})$
\end{enumerate}

\subsubsection*{Regression Loss Functions}

\begin{center}
\begin{tabular}{lll}
\toprule
Loss & $L(y, F)$ & Pseudo-residual $\tilde{r}$ \\
\midrule
Squared Error & $\frac{1}{2}(y - F)^2$ & $y - F$ \\
Absolute Error & $|y - F|$ & $\sign(y - F)$ \\
Huber ($\delta$) & see below & see below \\
Quantile ($\tau$) & see below & see below \\
\bottomrule
\end{tabular}
\end{center}

\textbf{Huber loss} (transition at $\delta$ set at the $\alpha$-quantile of $|r_i|$, default $\alpha = 0.9$):
\begin{equation}
L_\delta(r) = \begin{cases} \frac{1}{2}r^2 & |r| \leq \delta \\ \delta|r| - \frac{1}{2}\delta^2 & |r| > \delta \end{cases}
\end{equation}

\textbf{Quantile loss} (pinball loss, default $\tau = 0.5$):
\begin{equation}
L_\tau(r) = \begin{cases} \tau \cdot r & r \geq 0 \\ -(1-\tau) \cdot r & r < 0 \end{cases}
\end{equation}

For non-MSE losses, leaf values are overridden: median for absolute error, $\tau$-quantile for quantile loss, Huber-corrected values for Huber loss.

\subsubsection*{Binary Classification: Newton-Raphson Leaf Values}

For log-loss $L(y, F) = -[y\log\sigma(F) + (1-y)\log(1-\sigma(F))]$, the pseudo-residual is $\tilde{r}_i = y_i - \sigma(F(\mathbf{x}_i))$.

Instead of using the mean pseudo-residual as the leaf value, the implementation applies a Newton-Raphson step (matching sklearn):

\begin{equation}
\gamma_j = \frac{\sum_{i \in R_j} \tilde{r}_i}{\sum_{i \in R_j} p_i(1-p_i)}
\end{equation}

where $R_j$ is the set of samples in leaf $j$ and $p_i = \sigma(F(\mathbf{x}_i))$. This is derived from the second-order Taylor expansion of the loss around $F$, yielding a Newton step $-g/h$ where $g = \sum \tilde{r}_i$ and $h = \sum p_i(1-p_i)$.

\subsubsection*{Multiclass Classification}

For $K$ classes, maintain $K$ additive models $F_1, \ldots, F_K$ with softmax:
\begin{equation}
P(y=k|\mathbf{x}) = \frac{\exp(F_k(\mathbf{x}))}{\sum_{j=1}^K \exp(F_j(\mathbf{x}))}
\end{equation}

Each round fits $K$ separate regression trees. The pseudo-residual for class $k$ is:
\begin{equation}
\tilde{r}_{ik} = y_{ik} - P(y=k|\mathbf{x}_i)
\end{equation}
where $y_{ik} = \mathbb{1}[y_i = k]$.

Newton-Raphson leaf value with the $\frac{K-1}{K}$ scaling factor:
\begin{equation}
\gamma_{jk} = \frac{K-1}{K} \cdot \frac{\sum_{i \in R_{jk}} \tilde{r}_{ik}}{\sum_{i \in R_{jk}} |g_{ik}|(1 - |g_{ik}|)}
\end{equation}
where $g_{ik} = \tilde{r}_{ik}$. The $\frac{K-1}{K}$ factor corrects for the redundancy in the softmax parameterization.


%=============================================================================
\section{Support Vector Machines}
%=============================================================================

\subsection{Kernel SVM (SMO)}
\label{sec:ksvm}

\noindent\textit{Source:} \texttt{crates/scry-learn/src/svm/kernel.rs}

\subsubsection*{Primal Formulation}
\begin{equation}
\min_{\mathbf{w}, b, \boldsymbol{\xi}} \frac{1}{2}\|\mathbf{w}\|^2 + C\sum_{i=1}^n \xi_i \quad \text{s.t.} \quad y_i(\mathbf{w}^\top\phi(\mathbf{x}_i) + b) \geq 1 - \xi_i, \quad \xi_i \geq 0
\end{equation}

\subsubsection*{Dual Formulation}

Via Lagrange multipliers $\alpha_i$:
\begin{equation}
\max_{\boldsymbol{\alpha}} \sum_{i=1}^n \alpha_i - \frac{1}{2}\sum_{i,j}\alpha_i\alpha_j y_i y_j K(\mathbf{x}_i, \mathbf{x}_j) \quad \text{s.t.} \quad 0 \leq \alpha_i \leq C, \quad \sum_i \alpha_i y_i = 0
\end{equation}

\subsubsection*{KKT Conditions}

A solution $(\boldsymbol{\alpha}^*, b^*)$ satisfies:
\begin{align}
\alpha_i = 0 &\Rightarrow y_i f(\mathbf{x}_i) \geq 1 \\
0 < \alpha_i < C &\Rightarrow y_i f(\mathbf{x}_i) = 1 \quad \text{(support vector on margin)} \\
\alpha_i = C &\Rightarrow y_i f(\mathbf{x}_i) \leq 1 \quad \text{(inside margin or misclassified)}
\end{align}

\subsubsection*{Simplified SMO Algorithm}

At each step, select two multipliers $\alpha_i, \alpha_j$ violating the KKT conditions. The two-variable subproblem has an analytical solution.

\textbf{Bounds:}
\begin{equation}
L, H = \begin{cases}
(\max(0, \alpha_j - \alpha_i),\; \min(C, C + \alpha_j - \alpha_i)) & \text{if } y_i \neq y_j \\
(\max(0, \alpha_i + \alpha_j - C),\; \min(C, \alpha_i + \alpha_j)) & \text{if } y_i = y_j
\end{cases}
\end{equation}

\textbf{Update:}
\begin{align}
\eta &= 2K(\mathbf{x}_i, \mathbf{x}_j) - K(\mathbf{x}_i, \mathbf{x}_i) - K(\mathbf{x}_j, \mathbf{x}_j) \\
\alpha_j^{\text{new}} &= \alpha_j - \frac{y_j(E_i - E_j)}{\eta}, \quad \text{clipped to } [L, H]
\end{align}
where $E_i = f(\mathbf{x}_i) - y_i$ is the prediction error.

\subsubsection*{Kernel Functions}

\begin{center}
\begin{tabular}{ll}
\toprule
Kernel & $K(\mathbf{x}, \mathbf{x}')$ \\
\midrule
Linear & $\mathbf{x}^\top\mathbf{x}'$ \\
RBF & $\exp(-\gamma\|\mathbf{x} - \mathbf{x}'\|^2)$ \\
Polynomial & $(\mathbf{x}^\top\mathbf{x}' + c_0)^d$ \\
\bottomrule
\end{tabular}
\end{center}

\textbf{Gamma strategies:} Scale: $\gamma = \frac{1}{m \cdot \Var(\mathbf{X})}$ (sklearn default); Auto: $\gamma = \frac{1}{m}$.

\begin{theorem}[Mercer's Theorem (informal)]
A symmetric function $K: \mathcal{X} \times \mathcal{X} \to \mathbb{R}$ is a valid kernel (i.e., corresponds to an inner product in some feature space) if and only if the kernel matrix $[K(\mathbf{x}_i, \mathbf{x}_j)]_{ij}$ is positive semi-definite for all finite subsets $\{\mathbf{x}_1, \ldots, \mathbf{x}_n\}$.
\end{theorem}

\subsubsection*{Platt Scaling (Probability Calibration)}

Fits a logistic model $P(y=1|f) = \frac{1}{1+\exp(Af+B)}$ to the SVM decision values $f(\mathbf{x})$ using the training data, providing calibrated probability estimates.

%-----------------------------------------------------------------------------
\subsection{Linear SVM (Pegasos SGD)}
\label{sec:lsvm}

\noindent\textit{Source:} \texttt{crates/scry-learn/src/svm/linear.rs}

\subsubsection*{Primal Objective}
\begin{equation}
\min_\mathbf{w} \frac{\lambda}{2}\|\mathbf{w}\|^2 + \frac{1}{n}\sum_{i=1}^n \max(0, 1 - y_i\mathbf{w}^\top\mathbf{x}_i)
\end{equation}
where $\lambda = 1/(Cn)$.

\subsubsection*{Batch Sub-Gradient Update}
\begin{equation}
\mathbf{w}^{(t+1)} = \mathbf{w}^{(t)} - \eta_t\left[\lambda\mathbf{w}^{(t)} - \frac{1}{|\mathcal{B}|}\sum_{\substack{i \in \mathcal{B}:\\ y_i\mathbf{w}^\top\mathbf{x}_i < 1}} y_i\mathbf{x}_i\right]
\end{equation}

Learning rate: $\eta_t = \frac{1}{1 + 0.01t}$ (decaying).

\subsubsection*{Support Vector Regression (SVR)}

$\varepsilon$-insensitive loss:
\begin{equation}
L_\varepsilon(y, f) = \max(0, |y - f| - \varepsilon)
\end{equation}

Sub-gradient: $\nabla_\mathbf{w} = \lambda\mathbf{w} - \frac{1}{n}\sum_{i: |r_i| > \varepsilon} \sign(r_i)\mathbf{x}_i$ where $r_i = y_i - \mathbf{w}^\top\mathbf{x}_i$.

Learning rate: $\eta_t = \frac{1}{\lambda t}$.

\begin{proposition}[Pegasos Convergence]
Pegasos achieves an $\varepsilon$-accurate solution in $\tilde{O}(1/(\lambda\varepsilon))$ iterations, independent of the dataset size $n$ (Shalev-Shwartz et al., 2011).
\end{proposition}

\subsubsection*{Best-Weights Tracking}

The implementation tracks the weight vector achieving the lowest loss across all epochs and restores it at the end, providing robustness against late-training oscillations.


%=============================================================================
\section{Naive Bayes}
%=============================================================================

\subsection{Gaussian Naive Bayes}
\label{sec:gnb}

\noindent\textit{Source:} \texttt{crates/scry-learn/src/naive\_bayes/gaussian.rs}

\subsubsection*{Model}

Conditional independence assumption:
\begin{equation}
P(\mathbf{x}|y=k) = \prod_{j=1}^m P(x_j|y=k) = \prod_{j=1}^m \frac{1}{\sqrt{2\pi\sigma_{kj}^2}}\exp\left(-\frac{(x_j - \mu_{kj})^2}{2\sigma_{kj}^2}\right)
\end{equation}

\subsubsection*{MAP Classification}
\begin{equation}
\hat{y} = \argmax_k \left[\log\pi_k + \sum_{j=1}^m\left(-\frac{(x_j - \mu_{kj})^2}{2\sigma_{kj}^2} - \frac{1}{2}\log\sigma_{kj}^2 - \frac{1}{2}\log 2\pi\right)\right]
\end{equation}

\subsubsection*{Parameter Estimation}

Per-class statistics (with sample weights $w_i$):
\begin{equation}
\mu_{kj} = \frac{\sum_{i: y_i=k} w_i x_{ij}}{\sum_{i: y_i=k} w_i}, \quad \sigma_{kj}^2 = \frac{\sum_{i: y_i=k} w_i(x_{ij} - \mu_{kj})^2}{\sum_{i: y_i=k} w_i}, \quad \pi_k = \frac{\sum_{i: y_i=k} w_i}{\sum_i w_i}
\end{equation}

\subsubsection*{Variance Smoothing}

To prevent division by zero:
\begin{equation}
\sigma_{kj}^2 \leftarrow \sigma_{kj}^2 + \varepsilon, \quad \varepsilon = \texttt{var\_smoothing} \times \max_{k,j}(\sigma_{kj}^2)
\end{equation}
Default: $\texttt{var\_smoothing} = 10^{-9}$.

\subsubsection*{Numerical Stability: Log-Sum-Exp}
\begin{equation}
P(y=k|\mathbf{x}) = \frac{\exp(\ell_k - \ell_{\max})}{\sum_{j=1}^K \exp(\ell_j - \ell_{\max})}
\end{equation}
where $\ell_k = \log P(y=k) + \sum_j \log P(x_j|y=k)$ and $\ell_{\max} = \max_k \ell_k$.

%-----------------------------------------------------------------------------
\subsection{Multinomial Naive Bayes}

\noindent\textit{Source:} \texttt{crates/scry-learn/src/naive\_bayes/multinomial.rs}

\textbf{Likelihood:} $P(\mathbf{x}|y=k) \propto \prod_{j=1}^m \theta_{kj}^{x_j}$

\textbf{Parameter estimation with Laplace smoothing:}
\begin{equation}
\hat{\theta}_{kj} = \frac{N_{kj} + \alpha}{\sum_{j'=1}^m N_{kj'} + m\alpha}
\end{equation}
where $N_{kj} = \sum_{i: y_i=k} w_i x_{ij}$.

\textbf{Prediction:} $\log P(y=k|\mathbf{x}) \propto \log\pi_k + \sum_{j=1}^m x_j\log\hat{\theta}_{kj}$

%-----------------------------------------------------------------------------
\subsection{Bernoulli Naive Bayes}

\noindent\textit{Source:} \texttt{crates/scry-learn/src/naive\_bayes/bernoulli.rs}

\textbf{Likelihood} (binary features $x_j \in \{0,1\}$):
\begin{equation}
P(\mathbf{x}|y=k) = \prod_{j=1}^m p_{kj}^{x_j}(1-p_{kj})^{1-x_j}
\end{equation}

Unlike Multinomial NB, Bernoulli NB penalizes the \emph{absence} of features through the $(1-p_{kj})$ term.

\textbf{Parameters:}
\begin{equation}
\hat{p}_{kj} = \frac{N_{kj} + \alpha}{N_k + 2\alpha}
\end{equation}

\textbf{Prediction:}
\begin{equation}
\log P(y=k|\mathbf{x}) \propto \log\pi_k + \sum_{j=1}^m\left[x_j\log p_{kj} + (1-x_j)\log(1-p_{kj})\right]
\end{equation}

Optional binarization: $x_j \leftarrow \mathbb{1}[x_j > \text{threshold}]$ (default threshold: 0.0).


%=============================================================================
\section{K-Nearest Neighbors}
%=============================================================================

\noindent\textit{Source:} \texttt{crates/scry-learn/src/neighbors/knn.rs}

\subsection{Distance Metrics}

\begin{center}
\begin{tabular}{ll}
\toprule
Metric & $d(\mathbf{x}, \mathbf{x}')$ \\
\midrule
Euclidean & $\sqrt{\sum_j (x_j - x_j')^2}$ \\
Manhattan & $\sum_j |x_j - x_j'|$ \\
Cosine & $1 - \frac{\mathbf{x} \cdot \mathbf{x}'}{\|\mathbf{x}\|\|\mathbf{x}'\|}$ \\
\bottomrule
\end{tabular}
\end{center}

\begin{proposition}[Squared Euclidean Optimization]
For ranking nearest neighbors, squared Euclidean distance suffices since $\sqrt{\cdot}$ is monotonically increasing:
$d^2(\mathbf{x}, \mathbf{x}') < d^2(\mathbf{x}, \mathbf{x}'') \Leftrightarrow d(\mathbf{x}, \mathbf{x}') < d(\mathbf{x}, \mathbf{x}'')$.
This saves one \texttt{sqrt} per distance computation.
\end{proposition}

\subsection{Search Algorithms}

\textbf{Brute Force} ($O(nm)$ per query): Compute distances to all training points, use \texttt{select\_nth\_unstable} for $O(n)$ partial sort to find the $k$ nearest.

\textbf{KD-Tree} ($O(m\log n)$ average per query): Binary space partitioning. Restricted to Euclidean metric with $m \leq 20$ (KD-trees degrade in high dimensions due to the curse of dimensionality). Flat arena layout for cache-friendly traversal.

\textbf{Auto:} KD-tree when Euclidean and $m \leq 20$, brute force otherwise.

\subsection{Prediction}

\textbf{Classification:} $\hat{y} = \argmax_k \sum_{i \in \text{kNN}} w_i \cdot \mathbb{1}[y_i = k]$

\textbf{Regression:} $\hat{y} = \frac{\sum_{i \in \text{kNN}} w_i y_i}{\sum_{i \in \text{kNN}} w_i}$

where $w_i = 1$ (uniform) or $w_i = 1/d(\mathbf{x}_q, \mathbf{x}_i)$ (distance-weighted).


%=============================================================================
\section{Clustering}
%=============================================================================

\subsection{K-Means}
\label{sec:kmeans}

\noindent\textit{Source:} \texttt{crates/scry-learn/src/cluster/kmeans.rs}

\subsubsection*{Objective (WCSS / Inertia)}
\begin{equation}
\min_{\{C_k\}, \{\boldsymbol{\mu}_k\}} \sum_{k=1}^K \sum_{\mathbf{x}_i \in C_k} \|\mathbf{x}_i - \boldsymbol{\mu}_k\|^2
\end{equation}

\subsubsection*{Lloyd's Algorithm}
\begin{enumerate}[label=\arabic*., itemsep=0pt]
\item \textbf{Assign:} $c_i = \argmin_k \|\mathbf{x}_i - \boldsymbol{\mu}_k\|^2$
\item \textbf{Update:} $\boldsymbol{\mu}_k = \frac{1}{|C_k|}\sum_{\mathbf{x}_i \in C_k} \mathbf{x}_i$
\item Repeat until convergence.
\end{enumerate}

\begin{theorem}[K-Means Convergence]
Lloyd's algorithm converges in finitely many iterations.
\end{theorem}

\begin{proof}
(a) Each assignment step does not increase the objective (each point moves to its nearest centroid). (b) Each update step minimizes the objective for fixed assignments (the centroid is the mean, which minimizes squared distances). (c) The number of possible assignments is finite ($K^n$). Since the objective is non-increasing and bounded below by 0, and the number of distinct assignment states is finite, the algorithm must terminate.
\end{proof}

\begin{remark}
Lloyd's algorithm converges to a local minimum, not necessarily the global optimum.
\end{remark}

\subsubsection*{K-Means++ Initialization}

\begin{enumerate}[label=\arabic*., itemsep=0pt]
\item Choose $\boldsymbol{\mu}_1$ uniformly at random from $\{\mathbf{x}_i\}$.
\item For $k = 2, \ldots, K$: choose $\boldsymbol{\mu}_k = \mathbf{x}_i$ with probability $\propto D(\mathbf{x}_i)^2$ where $D(\mathbf{x}_i) = \min_{j<k}\|\mathbf{x}_i - \boldsymbol{\mu}_j\|$.
\end{enumerate}

\begin{theorem}[K-Means++ Approximation (Arthur \& Vassilvitskii, 2007)]
K-means++ achieves an expected cost at most $O(\log K)$ times the optimal cost: $\mathbb{E}[\text{cost}] \leq 8(\ln K + 2) \cdot \text{OPT}$.
\end{theorem}

The implementation runs \texttt{n\_init} independent runs (default: 10) and returns the result with lowest inertia.

%-----------------------------------------------------------------------------
\subsection{DBSCAN}
\label{sec:dbscan}

\noindent\textit{Source:} \texttt{crates/scry-learn/src/cluster/dbscan.rs}

\begin{definition}[Core Point]
A point $\mathbf{x}$ is a \emph{core point} if $|N_\varepsilon(\mathbf{x})| \geq \text{min\_samples}$, where $N_\varepsilon(\mathbf{x}) = \{\mathbf{x}' \in \mathcal{D} : d(\mathbf{x}, \mathbf{x}') \leq \varepsilon\}$.
\end{definition}

\begin{definition}[Density-Reachability]
$\mathbf{x}'$ is \emph{density-reachable} from $\mathbf{x}$ if there exists a chain of core points $\mathbf{x} = \mathbf{p}_1, \ldots, \mathbf{p}_l = \mathbf{x}'$ where each $\mathbf{p}_{i+1} \in N_\varepsilon(\mathbf{p}_i)$.
\end{definition}

A cluster is a maximal set of density-connected points. Points not reachable from any core point are classified as noise. The implementation uses a KD-tree for range queries ($O(\log n)$ average) when Euclidean and $m \leq 20$, brute force otherwise.

\subsubsection*{Properties}
\begin{itemize}[itemsep=0pt]
\item No cluster count required: discovers clusters automatically.
\item Finds non-convex clusters (unlike K-Means).
\item Core point assignments are deterministic; border point assignments may depend on processing order.
\end{itemize}

%-----------------------------------------------------------------------------
\subsection{Agglomerative Clustering}
\label{sec:agglo}

\noindent\textit{Source:} \texttt{crates/scry-learn/src/cluster/agglomerative.rs}

Bottom-up hierarchical clustering starting from $n$ singleton clusters, merging the closest pair until \texttt{n\_clusters} remain.

\subsubsection*{Linkage Criteria}

\begin{align}
d_{\text{single}}(A, B) &= \min_{\mathbf{a} \in A,\, \mathbf{b} \in B} d(\mathbf{a}, \mathbf{b}) \\
d_{\text{complete}}(A, B) &= \max_{\mathbf{a} \in A,\, \mathbf{b} \in B} d(\mathbf{a}, \mathbf{b}) \\
d_{\text{average}}(A, B) &= \frac{1}{|A||B|}\sum_{\mathbf{a} \in A}\sum_{\mathbf{b} \in B} d(\mathbf{a}, \mathbf{b}) \\
d_{\text{Ward}}(A, B) &= \frac{|A||B|}{|A|+|B|}\|\boldsymbol{\mu}_A - \boldsymbol{\mu}_B\|^2
\end{align}

Ward minimizes total within-cluster variance after merging.

\subsubsection*{Lance-Williams Recurrence}

Distances from a merged cluster $A \cup B$ to cluster $C$:
\begin{align}
d_{\text{Ward}}(A \cup B, C) &= \frac{(|A|+|C|)d(A,C) + (|B|+|C|)d(B,C) - |C|\cdot d(A,B)}{|A|+|B|+|C|} \\
d_{\text{single}}(A \cup B, C) &= \min(d(A,C), d(B,C)) \\
d_{\text{complete}}(A \cup B, C) &= \max(d(A,C), d(B,C)) \\
d_{\text{average}}(A \cup B, C) &= \frac{|A|\cdot d(A,C) + |B|\cdot d(B,C)}{|A|+|B|}
\end{align}

The implementation uses a \texttt{BinaryHeap} (max-heap with negated distances) for $O(\log n)$ minimum extraction per merge.


%=============================================================================
\section{Neural Networks}
%=============================================================================

\noindent\textit{Source:} \texttt{crates/scry-learn/src/neural/network.rs}

\subsection{Dense Layer Forward Pass}

For layer $l$ with weights $\mathbf{W}^{(l)} \in \mathbb{R}^{d_{\mathrm{out}} \times d_{\mathrm{in}}}$, biases $\mathbf{b}^{(l)}$, and activation $\sigma$:
\begin{align}
\mathbf{z}^{(l)} &= \mathbf{W}^{(l)}\mathbf{a}^{(l-1)} + \mathbf{b}^{(l)} \\
\mathbf{a}^{(l)} &= \sigma(\mathbf{z}^{(l)})
\end{align}

\subsection{Activation Functions}

\begin{center}
\begin{tabular}{lll}
\toprule
Activation & $\sigma(z)$ & $\sigma'(z)$ \\
\midrule
ReLU & $\max(0, z)$ & $\mathbb{1}[z > 0]$ \\
Tanh & $\tanh(z)$ & $1 - \tanh^2(z)$ \\
Sigmoid & $\frac{1}{1+e^{-z}}$ & $\sigma(z)(1-\sigma(z))$ \\
Identity & $z$ & $1$ \\
\bottomrule
\end{tabular}
\end{center}

\subsection{Backpropagation}

For a network with $L$ layers, given loss gradient $\frac{\partial\mathcal{L}}{\partial\mathbf{a}^{(L)}}$:

\textbf{Output layer:}
\begin{equation}
\boldsymbol{\delta}^{(L)} = \frac{\partial\mathcal{L}}{\partial\mathbf{a}^{(L)}} \odot \sigma'(\mathbf{z}^{(L)})
\end{equation}

\textbf{Hidden layers} ($l = L-1, \ldots, 1$):
\begin{equation}
\boldsymbol{\delta}^{(l)} = (\mathbf{W}^{(l+1)})^\top\boldsymbol{\delta}^{(l+1)} \odot \sigma'(\mathbf{z}^{(l)})
\end{equation}

\textbf{Parameter gradients:}
\begin{align}
\frac{\partial\mathcal{L}}{\partial\mathbf{W}^{(l)}} &= \boldsymbol{\delta}^{(l)}(\mathbf{a}^{(l-1)})^\top + \alpha\mathbf{W}^{(l)} \\
\frac{\partial\mathcal{L}}{\partial\mathbf{b}^{(l)}} &= \boldsymbol{\delta}^{(l)}
\end{align}
where $\alpha\mathbf{W}^{(l)}$ is the L2 regularization gradient (weights only, not biases).

\subsection{Loss Functions}

\textbf{Cross-entropy} (classification):
\begin{equation}
\mathcal{L} = -\frac{1}{n}\sum_{i=1}^n \log\frac{e^{z_{i,y_i}}}{\sum_k e^{z_{ik}}}, \quad \frac{\partial\mathcal{L}}{\partial z_{ik}} = \softmax(z_{ik}) - \mathbb{1}[k = y_i]
\end{equation}

\textbf{MSE} (regression):
\begin{equation}
\mathcal{L} = \frac{1}{n}\sum_{i=1}^n(f(\mathbf{x}_i) - y_i)^2, \quad \frac{\partial\mathcal{L}}{\partial f_i} = \frac{2}{n}(f_i - y_i)
\end{equation}

\subsection{Weight Initialization}

\textbf{He} (for ReLU): $W_{ij} \sim \mathcal{N}(0, 2/d_{\mathrm{in}})$

\textbf{Xavier/Glorot} (for Sigmoid/Tanh): $W_{ij} \sim \mathcal{N}(0, 2/(d_{\mathrm{in}} + d_{\mathrm{out}}))$

\subsection{Optimizers}

\textbf{Adam} (default): $\beta_1 = 0.9$, $\beta_2 = 0.999$, $\epsilon = 10^{-8}$.
\begin{align}
\mathbf{m}_t &= \beta_1\mathbf{m}_{t-1} + (1-\beta_1)\mathbf{g}_t \\
\mathbf{v}_t &= \beta_2\mathbf{v}_{t-1} + (1-\beta_2)\mathbf{g}_t^2 \\
\hat{\mathbf{m}}_t &= \mathbf{m}_t/(1-\beta_1^t), \quad \hat{\mathbf{v}}_t = \mathbf{v}_t/(1-\beta_2^t) \\
\boldsymbol{\theta}_t &= \boldsymbol{\theta}_{t-1} - \eta\hat{\mathbf{m}}_t/(\sqrt{\hat{\mathbf{v}}_t} + \epsilon)
\end{align}

\textbf{SGD with Nesterov momentum} ($\mu = 0.9$):
\begin{align}
\mathbf{v}_t &= \mu\mathbf{v}_{t-1} - \eta\nabla f(\boldsymbol{\theta}_{t-1} + \mu\mathbf{v}_{t-1}) \\
\boldsymbol{\theta}_t &= \boldsymbol{\theta}_{t-1} + \mathbf{v}_t
\end{align}

\subsection{Conv2D}

\noindent\textit{Source:} \texttt{crates/scry-learn/src/neural/conv.rs}

2D convolution via \textbf{im2col + GEMM}. Input: $[\text{batch}, C_{\text{in}}, H, W]$ (channels-first).

Output dimensions:
\begin{equation}
H_{\text{out}} = \frac{H + 2p - k}{s} + 1, \quad W_{\text{out}} = \frac{W + 2p - k}{s} + 1
\end{equation}
where $k$ is kernel size, $s$ is stride, $p$ is padding.


%=============================================================================
\section{Anomaly Detection}
%=============================================================================

\subsection{Isolation Forest}
\label{sec:iforest}

\noindent\textit{Source:} \texttt{crates/scry-learn/src/anomaly/iforest.rs}

\subsubsection*{Intuition}

Anomalies are ``few and different''---they can be isolated with fewer random partitions than normal points, yielding shorter average path lengths in random trees.

\subsubsection*{Isolation Tree Construction}

For a subsample $\mathcal{S}$ of size $\psi$:
\begin{enumerate}[label=\arabic*., itemsep=0pt]
\item If $|\mathcal{S}| \leq 1$ or depth $\geq$ max\_depth: return leaf with size $|\mathcal{S}|$.
\item Select random feature $j$ and random threshold $\theta \in [\min_j, \max_j]$.
\item Partition: $\mathcal{S}_L = \{x : x_j \leq \theta\}$, $\mathcal{S}_R = \{x : x_j > \theta\}$.
\item Recurse on $\mathcal{S}_L$ and $\mathcal{S}_R$.
\end{enumerate}

\subsubsection*{Path Length and Anomaly Score}

The path length includes a correction for the leaf's external path length:
\begin{equation}
h(\mathbf{x}) = \text{depth} + c(|\text{leaf}|)
\end{equation}
where:
\begin{equation}
c(n) = \begin{cases} 2H(n-1) - \frac{2(n-1)}{n} & n > 2 \\ 1 & n = 2 \\ 0 & n \leq 1 \end{cases}
\end{equation}
$H(k) = \ln(k) + \gamma_E \approx \ln(k) + 0.5772$ is the harmonic number (Euler-Mascheroni constant).

\textbf{Anomaly score:}
\begin{equation}
s(\mathbf{x}, \psi) = 2^{-E[h(\mathbf{x})]/c(\psi)}
\end{equation}

\textbf{Interpretation:} $s \approx 1$: anomaly (short paths); $s \approx 0.5$: normal; $s \approx 0$: extremely dense region.


%=============================================================================
\section{Preprocessing}
%=============================================================================

\subsection{PCA (Principal Component Analysis)}
\label{sec:pca}

\noindent\textit{Source:} \texttt{crates/scry-learn/src/preprocess/pca.rs}

\subsubsection*{Objective}
\begin{equation}
\mathbf{w}_1 = \argmax_{\|\mathbf{w}\|=1} \Var(\mathbf{X}\mathbf{w}) = \argmax_{\|\mathbf{w}\|=1} \mathbf{w}^\top\boldsymbol{\Sigma}\mathbf{w}
\end{equation}

\subsubsection*{Covariance Matrix}
\begin{equation}
\boldsymbol{\Sigma} = \frac{1}{n-1}\sum_{i=1}^n (\mathbf{x}_i - \bar{\mathbf{x}})(\mathbf{x}_i - \bar{\mathbf{x}})^\top
\end{equation}

\subsubsection*{Jacobi Eigenvalue Algorithm}

The implementation uses Jacobi rotation to diagonalize $\boldsymbol{\Sigma}$:

\begin{enumerate}[label=\arabic*., itemsep=0pt]
\item Find off-diagonal element with largest absolute value: $|\Sigma_{pq}|$.
\item Compute rotation angle: $\tan 2\theta = \frac{2\Sigma_{pq}}{\Sigma_{pp} - \Sigma_{qq}}$.
\item Apply the $(p,q)$-rotation, zeroing out $\Sigma_{pq}$ and $\Sigma_{qp}$.
\item Repeat sweeps until $\|\mathrm{off\text{-}diag}(\boldsymbol{\Sigma})\|_F < 10^{-12}$.
\end{enumerate}

Diagonal elements converge to eigenvalues; the accumulated rotation matrix gives eigenvectors.

\begin{proposition}[Convergence]
Each Jacobi sweep reduces the off-diagonal Frobenius norm. The method converges quadratically for symmetric matrices (Golub \& Van Loan, 2013).
\end{proposition}

\subsubsection*{Transform and Whitening}

\textbf{Transform:} $\mathbf{z}_i = \mathbf{V}_k^\top(\mathbf{x}_i - \bar{\mathbf{x}})$

\textbf{Whitening:} $\tilde{\mathbf{z}}_i = \boldsymbol{\Lambda}_k^{-1/2}\mathbf{V}_k^\top(\mathbf{x}_i - \bar{\mathbf{x}})$

\textbf{Inverse:} $\hat{\mathbf{x}}_i = \mathbf{V}_k\mathbf{z}_i + \bar{\mathbf{x}}$ (lossy when $k < m$)

%-----------------------------------------------------------------------------
\subsection{Scalers}

\textbf{Standard:} $x'_j = \frac{x_j - \mu_j}{\sigma_j}$ \quad (zero-variance features unchanged)

\textbf{MinMax:} $x'_j = \frac{x_j - \min_j}{\max_j - \min_j}$ \quad (maps to $[0,1]$)

\textbf{Robust:} $x'_j = \frac{x_j - \text{median}_j}{\text{IQR}_j}$ \quad (robust to outliers)

\begin{remark}
Robust Scaler uses median and IQR, which are unaffected by extreme values. A single outlier can shift the mean and inflate the standard deviation, but has minimal effect on the median and interquartile range.
\end{remark}


%=============================================================================
\section{Model Selection}
%=============================================================================

\subsection{Grid Search CV}

\noindent\textit{Source:} \texttt{crates/scry-learn/src/search/grid.rs}

Exhaustive search over the Cartesian product of parameter values, evaluated via $k$-fold cross-validation. If $p_1$ has $v_1$ values and $p_2$ has $v_2$ values, evaluates all $v_1 \times v_2$ combinations.

Complexity: $O(|\text{grid}| \times k \times T_{\text{fit}})$.

\subsection{Randomized Search CV}

\noindent\textit{Source:} \texttt{crates/scry-learn/src/search/random.rs}

Samples \texttt{n\_iter} random combinations from the grid (Fisher-Yates shuffle, then take first \texttt{n\_iter}).

\begin{proposition}[Random vs Grid (Bergstra \& Bengio, 2012)]
When only a few hyperparameters matter, random search is more efficient than grid search because it explores more distinct values per important parameter.
\end{proposition}

\subsection{Bayesian Search CV (TPE)}

\noindent\textit{Source:} \texttt{crates/scry-learn/src/search/bayes.rs}

Tree-Structured Parzen Estimator (Bergstra et al., 2011). After an initial random exploration phase ($n_{\text{initial}}$ iterations):

\begin{enumerate}[label=\arabic*., itemsep=0pt]
\item Split observed results at the $\gamma$-quantile (default $\gamma = 0.25$) into ``good'' ($\ell$) and ``bad'' ($g$) groups.
\item Build factored 1D kernel density estimates for each group.
\item Draw 100 random candidates and select the one maximizing:
\begin{equation}
\text{EI}(\mathbf{x}) \propto \frac{\ell(\mathbf{x})}{g(\mathbf{x})}
\end{equation}
\end{enumerate}

\subsubsection*{Kernel Density Estimation}

\textbf{Continuous parameters:} 1D Gaussian KDE with Scott's rule bandwidth $h = n^{-1/5}$, parameters normalized to $[0,1]$.

\textbf{Categorical parameters:} Frequency-based with Laplace smoothing:
\begin{equation}
P(x = v_i) = \frac{\text{count}(v_i) + 1}{\sum_j(\text{count}(v_j) + 1)}
\end{equation}


%=============================================================================
\section*{References}
%=============================================================================

\begin{enumerate}[label={[\arabic*]}, itemsep=2pt]
\item Nocedal, J. \& Wright, S.J. (2006). \textit{Numerical Optimization}, 2nd ed. Springer.
\item Breiman, L. et al. (1984). \textit{Classification and Regression Trees}. Wadsworth.
\item Breiman, L. (2001). Random Forests. \textit{Machine Learning}, 45(1), 5--32.
\item Friedman, J.H. (2001). Greedy function approximation: A gradient boosting machine. \textit{Annals of Statistics}, 29(5), 1189--1232.
\item Platt, J.C. (1998). Sequential Minimal Optimization. Microsoft Research TR MSR-TR-98-14.
\item Shalev-Shwartz, S. et al. (2011). Pegasos: Primal Estimated sub-GrAdient SOlver for SVM. \textit{Mathematical Programming}, 127(1), 3--30.
\item Ester, M. et al. (1996). A density-based algorithm for discovering clusters in large spatial databases with noise. \textit{KDD}, 226--231.
\item Arthur, D. \& Vassilvitskii, S. (2007). k-means++: The advantages of careful seeding. \textit{SODA}, 1027--1035.
\item Liu, F.T. et al. (2008). Isolation Forest. \textit{ICDM}, 413--422.
\item Tseng, P. (2001). Convergence of a block coordinate descent method for nondifferentiable minimization. \textit{JOTA}, 109(3), 475--494.
\item Bergstra, J. et al. (2011). Algorithms for Hyper-Parameter Optimization. \textit{NeurIPS}.
\item Bergstra, J. \& Bengio, Y. (2012). Random Search for Hyper-Parameter Optimization. \textit{JMLR}, 13, 281--305.
\item Golub, G.H. \& Van Loan, C.F. (2013). \textit{Matrix Computations}, 4th ed. Johns Hopkins University Press.
\end{enumerate}

\end{document}
